\subsection{Change-point model: conditional transport RJMCMC with collapsed rates}

We consider the classical coal-mining-disaster change-point benchmark. Let
$y=\{t_1,\dots,t_n\}\subset[0,L]$ denote the ordered event times, with
$n=190$ and $L=40907$ in the data set used here. For a model with
$k\in\{0,\dots,k_{\max}\}$ change points, the observation window is partitioned
into $k+1$ segments with lengths $\ell_1,\dots,\ell_{k+1}$ satisfying
$\ell_i>0$ and $\sum_{i=1}^{k+1}\ell_i=L$. Writing
\[
w_i=\ell_i/L,\qquad w=(w_1,\dots,w_{k+1})\in\Delta_k,
\]
the associated change-point locations are
\[
\tau_j=L\sum_{i=1}^j w_i,\qquad j=1,\dots,k.
\]
Conditional on $(k,w,h)$, where $h=(h_1,\dots,h_{k+1})$ are the segmentwise
Poisson intensities, the data are modeled as a piecewise homogeneous Poisson
process, so that
\[
p(y\mid k,w,h)\propto \prod_{i=1}^{k+1} h_i^{N_i}\exp(-\ell_i h_i),
\]
where $N_i$ is the number of events in segment $i$.

We assign a Poisson prior on $k$, truncated to $\{0,\dots,k_{\max}\}$,
\[
k \sim \text{Poisson}(\lambda),\qquad \lambda=3,
\]
a Dirichlet prior on the normalized segment lengths,
\[
w\mid k \sim \operatorname{Dirichlet}(2,\dots,2),
\]
and conditionally independent Gamma priors on the intensities,
\[
h_i\mid k \stackrel{\text{i.i.d.}}{\sim} \Gamma(\alpha,\beta),
\qquad \alpha=1,\quad \beta=200,
\]
under the shape-rate parameterization. In the version studied here, the rate
vector $h$ is integrated out analytically. Hence the RJMCMC state is reduced to
$(k,w)$, or equivalently to $(k,s)$ after reparameterization below. The
resulting marginal likelihood is
\[
p(y\mid k,w)\propto
\prod_{i=1}^{k+1}
\frac{\Gamma(\alpha+N_i)}{\Gamma(\alpha)}
\frac{\beta^\alpha}{(\beta+\ell_i)^{\alpha+N_i}},
\]
so that the collapsed posterior is
\[
\pi(k,w\mid y)\propto p(k)\,p(w\mid k)\,p(y\mid k,w).
\]

To work in Euclidean space, we use a stick-breaking parameterization of the
simplex. For model $k$, let $s=(s_1,\dots,s_k)\in\mathbb{R}^k$ and define
\[
z_i=\sigma\!\bigl(s_i-\log(k+1-i)\bigr),\qquad i=1,\dots,k,
\]
where $\sigma(x)=(1+e^{-x})^{-1}$. The simplex coordinates are then recovered
through
\[
w_1=z_1,\qquad
w_i=z_i\prod_{j=1}^{i-1}(1-z_j),\ \ i=2,\dots,k,\qquad
w_{k+1}=\prod_{j=1}^{k}(1-z_j).
\]
This yields a $k$-dimensional unconstrained representation for model $k$. The
target density on $s$ is obtained by combining the truncated Poisson prior on
$k$, the Dirichlet density on $w(s)$, the Jacobian of the stick-breaking map,
and the collapsed likelihood above.

For trans-dimensional proposals we use a conditional transport construction.
Rather than learning separate maps for each neighboring pair of models, we fit a
single conditional normalizing flow on a padded target of fixed dimension
$k_{\max}$. For a state in model $k$, the first $k$ coordinates correspond to
the active unconstrained spacing variables, while the remaining $k_{\max}-k$
coordinates are filled with independent Gaussian reference variables with
standard deviation $\sigma_u=1$. A binary context vector indicating the active
coordinates is supplied to the transport map. In this way, all models are
embedded in the same ambient space and the transport is conditioned only on the
current model dimension.

Given a current state $(k,s)$, let $x^{(k)}\in\mathbb{R}^{k_{\max}}$ denote the
corresponding padded representation. The conditional flow maps $x^{(k)}$ to a
base-space variable $z\in\mathbb{R}^{k_{\max}}$. To propose a jump to a
neighboring model $k'\in\{k-1,k+1\}$, we keep the same base-space point $z$ and
decode it under the context associated with model $k'$. This produces a padded
proposal $x^{(k')}$ whose first $k'$ coordinates determine the proposed
unconstrained spacing vector $s'$. Since the inactive coordinates are modeled
explicitly through the Gaussian reference law, the Metropolis--Hastings ratio
contains the corresponding reference-density contribution together with the
forward and reverse Jacobian terms induced by the transport map. For interior
models the proposal chooses $k-1$ and $k+1$ with equal probability, whereas at
the boundaries the unique admissible neighboring model is selected.

Within-model updates are performed with probability
$p_{\mathrm{within}}=0.5$. In that case we apply a Gaussian random walk to the
active unconstrained coordinates, using a covariance estimated from calibration
particles and multiplied by a scale factor of $0.35$. With the remaining
probability, a birth/death proposal is generated by the conditional transport
mechanism described above. The resulting Markov kernel therefore combines local
within-model exploration with global trans-dimensional transport moves.

In the numerical study we use $k_{\max}=10$. The conditional flow is trained by
variational inference with annealing for $10{,}000$ gradient iterations, using
$256$ Monte Carlo samples per update. For dimensions larger than one we employ a
conditional RealNVP architecture with a trainable base distribution. After the
transport has been fitted, RJMCMC is run for either $40{,}000$ or $100{,}000$
iterations, depending on the experiment. The chain is initialized from a small
set of flow-generated candidates by selecting the one with the largest posterior
score under the collapsed target.

Posterior model probabilities are estimated from the empirical occupancy of the
sampled model index. Conditional on the modal model, posterior summaries for the
segment intensities are obtained from the collapsed conditional mean
\[
\mathrm{E}[h_i\mid y,k,w]=\frac{\alpha+N_i}{\beta+\ell_i}.
\]
This formulation separates the effect of the transport proposal from the
additional Monte Carlo variability that would arise from explicitly sampling the
nuisance rate parameters.
